\begin{ouverture}[courbe]{6}{Proportionnalité, ratios et pourcentages}{Partie 2 · Grandeurs et fonctions}
{Coefficients · Échelles · Évolutions · Comparaisons}
\objectif{Reconnaître une proportionnalité.}
\objectif{Utiliser un ratio ou une échelle.}
\objectif{Appliquer et interpréter des pourcentages.}
\prerequis{Fractions et produits.}
\end{ouverture}
\pagestyle{otMain}\essentiel\fontsize{11}{15}\selectfont
\essSection{1}{Reconnaître une proportionnalité}
\begin{definition}Deux grandeurs sont proportionnelles si l’on obtient les valeurs de l’une en multipliant celles de l’autre par un même coefficient. Le graphique est alors une droite passant par l’origine.\end{definition}
\begin{savoirfaire}[Choisir une méthode]Utiliser le passage à l’unité, un coefficient, la linéarité ou un produit en croix. Si 3 kg coûtent 7,20 €, 1 kg coûte 2,40 €, et 5 kg coûtent 12 €.\end{savoirfaire}
\begin{definition}Un ratio $a:b$ partage une quantité en $a+b$ parts égales, avec $a$ parts pour le premier terme. Pour un ratio $2:3$ sur 35 L, les volumes sont $14$ L et $21$ L.\end{definition}
\begin{definition}Une échelle est le rapport longueur sur le plan sur longueur réelle, dans la même unité. À l’échelle $1/500$, 1 cm sur le plan représente 500 cm, soit 5 m.\end{definition}
\begin{attention}Un tarif avec un forfait fixe n’est pas proportionnel au nombre d’unités. Deux valeurs qui augmentent ensemble ne sont pas forcément proportionnelles.\end{attention}
\clearpage
\essSection{2}{Appliquer et comprendre un pourcentage}
\begin{propriete}Calculer $t\,\%$ d’une quantité revient à la multiplier par $t/100$. Une hausse de $t\,\%$ multiplie par $1+t/100$ ; une baisse par $1-t/100$.\end{propriete}
\begin{exemple}Après une baisse de 15\,\%, un prix de 80 € devient $80\times0{,}85=68$ €. Retrouver le prix initial exige de diviser par 0,85.\end{exemple}
\begin{savoirfaire}[Évolutions successives]Multiplier les coefficients, sans additionner les taux. Une hausse de 20\,\% suivie d’une baisse de 20\,\% donne $1{,}2\times0{,}8=0{,}96$ : la baisse totale est de 4\,\%.\end{savoirfaire}
\begin{propriete}Le taux d’évolution de $V_i>0$ à $V_f$ vaut $(V_f-V_i)/V_i$. Préciser la valeur initiale choisie comme référence.\end{propriete}
\begin{attention}Passer de 30\,\% à 36\,\% représente 6 points de pourcentage, mais une augmentation relative de $6/30=20\,\%$. Ces deux informations ne répondent pas à la même question.\end{attention}
\clearpage\rubrique{Découverte}
\begin{blocLivre}{ex-6-1}\begin{exercice}[1]{Passer par l’unité}
Quatre cahiers coûtent 10,80 €. Le prix est proportionnel au nombre de cahiers. Quel prix pour sept cahiers ?
\end{exercice}\end{blocLivre}
\begin{blocLivre}{ex-6-2}\begin{exercice}[1]{Un tableau proportionnel ?}
Les couples quantité-prix sont $(2;7)$, $(5;17{,}5)$ et $(8;28)$. Le tableau est-il proportionnel ?
\end{exercice}\end{blocLivre}
\begin{blocLivre}{ex-6-3}\begin{exercice}[1]{Appliquer une réduction}
Un sac coûte 64 €. Son prix baisse de 25 \%. Calculer le nouveau prix.
\end{exercice}\end{blocLivre}
\begin{blocLivre}{ex-6-4}\begin{exercice}[1]{Calculer un taux}
Un effectif passe de 160 à 184. Quel est son taux d’augmentation ?
\end{exercice}\end{blocLivre}
\begin{blocLivre}{ex-6-5}\begin{exercice}[1]{Partager selon un ratio}
Partager 56 billes selon le ratio $3:5$.
\end{exercice}\end{blocLivre}

\clearpage\rubrique{Entraînement}
\begin{blocLivre}{ex-6-6}\begin{exercice}[2]{Retrouver un prix initial}
Après une hausse de 8 \%, un abonnement coûte 81 €. Quel était son prix ?
\end{exercice}\end{blocLivre}
\begin{blocLivre}{ex-6-7}\begin{exercice}[2]{Deux baisses}
Un prix de 250 € baisse de 10 \%, puis de 20 \%. Calculer le prix final et la réduction globale.
\end{exercice}\end{blocLivre}
\begin{blocLivre}{ex-6-8}\begin{exercice}[2]{Une échelle}
Deux lieux distants de 7,2 km sont séparés de 9 cm sur une carte. Déterminer l’échelle.
\end{exercice}\end{blocLivre}
\begin{blocLivre}{ex-6-9}\begin{exercice}[2]{Une recette à trois termes}
Une boisson est préparée selon le ratio sirop : eau : jus $=1:6:3$. Pour 4,5 L de boisson, calculer chaque volume.
\end{exercice}\end{blocLivre}
\begin{blocLivre}{ex-6-10}\begin{exercice}[2]{Un tarif avec forfait}
Une livraison coûte 5 € de préparation, puis 2 € par colis. Le coût est-il proportionnel au nombre de colis ? Calculer pour 3 et 6 colis.
\end{exercice}\end{blocLivre}

\clearpage\rubrique{Problèmes}
\begin{blocLivre}{ex-6-11}\begin{exercice}[3]{Comparer deux promotions}
Un article coûte 120 €. Offre A : réduction de 30 €. Offre B : réduction de 20 \%, puis bon d’achat de 10 € utilisable immédiatement sur cet achat. Comparer les paiements.
\end{exercice}\end{blocLivre}
\begin{blocLivre}{ex-6-12}\begin{exercice}[3]{Un mélange à corriger}
On mélange 2 L de sirop avec 8 L d’eau. On veut que le sirop représente 16 \% du mélange sans en retirer. Quel volume d’eau ajouter ?
\end{exercice}\end{blocLivre}
\begin{blocLivre}{ex-6-13}\begin{exercice}[3]{Deux pourcentages emboîtés}
Dans un collège de 720 élèves, 45 \% sont inscrits à l’association sportive. Parmi les inscrits, $2/9$ pratiquent le badminton. Quel est leur effectif et leur part dans le collège ?
\end{exercice}\end{blocLivre}
\begin{blocLivre}{ex-6-14}\begin{exercice}[3]{Revenir après une baisse}
Un prix baisse de 20 \%. Quelle augmentation en pourcentage permet de revenir au prix initial ?
\end{exercice}\end{blocLivre}
\begin{blocLivre}{ex-6-15}\begin{exercice}[3]{Lire une annonce}
La part des élèves venant à vélo passe de 12 \% à 15 \%. Une annonce parle de « 3 \% de hausse ». Expliquer précisément l’évolution. Sur 800 élèves à effectif constant, combien de cyclistes supplémentaires ?
\end{exercice}\end{blocLivre}

\clearpage
